\documentclass[11pt]{article}

% --- Préambule compatible arXiv (paquets standard uniquement) ---
\usepackage[utf8]{inputenc}
\usepackage[T1]{fontenc}
\usepackage[french]{babel}
\usepackage{amsmath,amssymb,amsthm}
\usepackage{booktabs}
\usepackage{geometry}
\usepackage{hyperref}
\usepackage{xcolor}
\geometry{margin=1in}
\hypersetup{colorlinks=true,linkcolor=blue,citecolor=blue,urlcolor=blue}

\newtheorem{definition}{Définition}
\newtheorem{proposition}{Proposition}
\newtheorem{theorem}{Théorème}

\newcommand{\dcausal}{d_{\mathrm{c}}}
\newcommand{\Nk}{\mathcal{N}_k}
\newcommand{\code}[1]{\texttt{#1}}

\title{\textbf{La MMU cognitive :\\ une mémoire de travail déterministe et auditable pour les agents de code à base de LLM\\ \large (et un compte rendu honnête de pourquoi la recherche causale ne bat pas le RAG)}}
\author{Le projet CCOS\thanks{Causal Context Operating System (CCOS), un prototype
de recherche ouvert. Le code source, les scripts de reproduction et cet article :
\url{https://github.com/CHECKUPAUTO/CCOS}.}}
\date{\today}

\begin{document}
\maketitle

\begin{abstract}
La mémoire de travail d'un agent de code à base de LLM est gérée de façon opaque :
une pile de récupération (RAG, GraphRAG, frameworks d'agents) pagine un dépôt dans
une fenêtre bornée et mute cette fenêtre de manière probabiliste au fil d'une
session ; aussi, lorsqu'un agent dérive après des dizaines d'étapes, il est
impossible de dire \emph{pourquoi} ni \emph{quand} sa représentation du projet
s'est corrompue. CCOS traite la mémoire de l'agent comme un système d'exploitation
traite la RAM : le code source est analysé en un graphe causal \emph{event-sourced},
les nœuds sont scorés et paginés sous un budget de jetons, et chaque transition
d'état est enregistrée dans un journal chaîné par hachage qui se rejoue bit pour
bit. Nous donnons une définition précise et déterministe d'une \emph{région causale}
(une composante connexe du graphe des liens causaux inter-fichiers) et de la
distance causale (un plus court chemin pondéré), et prouvons que les régions et la
fenêtre paginée sont une fonction pure du graphe, se reconstruisant exactement au
rejeu --- la propriété qui rend la mémoire d'un agent \emph{auditable} et
\emph{déboguable par voyage dans le temps}. Nous demandons ensuite si cette structure
causale fait de CCOS un \emph{meilleur récupérateur} qu'une référence de base, et
rapportons --- honnêtement --- que non. Mesurée sur $70$ commits réels de correction
de bogues issus de trois bibliothèques (« crates ») matures, la sélection causale
\textbf{fait jeu égal avec un simple récupérateur lexical TF-IDF} et lui est
inférieure à budget serré ; sur des bogues réels, les fichiers d'un correctif
partagent leur vocabulaire, si bien que la similarité lexicale les trouve aussi, et
une preuve de concept amorçant CCOS depuis une trace de crash est elle aussi battue
par un RAG-sur-le-message-d'erreur (les messages d'erreur de Rust nomment la cause).
Une simulation sans LLM \emph{sépare} bien les deux lorsqu'une cause est construite
pour être lexicalement dissemblable de son symptôme, et un premier essai embarqué
avec un modèle de $7$ milliards de paramètres le reproduit --- mais ce régime ne se
produit pas sur du vrai code. La contribution n'est donc \emph{pas} une meilleure
récupération. C'est l'infrastructure qui manque structurellement à une pile de
récupération : une mémoire d'agent \emph{déterministe, rejouable, auditable} dans
laquelle on peut rembobiner l'état exact du contexte à n'importe quelle étape et le
rejouer sous d'autres paramètres (\emph{débogage par voyage dans le temps}). À notre
connaissance, CCOS est le \emph{premier} système à traiter la \emph{mémoire de travail
elle-même} d'un agent comme un sous-système transactionnel --- déterministe, chaîné
par hachage, rejouable bit-à-bit et \emph{déboguable a posteriori} : un enregistreur
de vol de l'attention, avec un point d'arrêt sur l'étape exacte où la bonne
information a été évincée de la fenêtre. Nous
publions le moteur, le harnais de validation honnête et cet article, en étiquetant
partout ce qui est \emph{prouvé} (structure formelle, déterminisme, rejeu),
\emph{mesuré} (couverture de récupération --- où CCOS ne bat \emph{pas} le RAG
lexical sur des bogues réels), et le seul axe où CCOS \emph{gagne} :
\textbf{l'efficacité}. Exécuté de bout en bout sur de vraies corrections
mono-fichier avec un modèle de 30B et une nouvelle tentative avec le compilateur
dans la boucle, CCOS et le RAG résolvent \emph{à égalité} ($2/10$), mais CCOS le
fait sur $776$ jetons de contexte contre $5366$ pour le récupérateur lexical --- une
réduction de $6{,}9\times$ (et de $4{,}1$ à $9{,}1\times$ sur $51$ correctifs issus
de trois crates, sans modèle), car la région causale s'arrête au working set au lieu
de remplir un top-$k$ jusqu'au budget. CCOS n'est donc pas un meilleur récupérateur
ou solveur, mais un récupérateur plus \emph{frugal} et \emph{auditable} --- il
s'auto-calibre, sans $k$ à régler. (Sur sa propre arborescence source, la sélection
par région couvre le voisinage à $k$ sauts d'une tâche avec un rappel de $0{,}97$
contre $0{,}35$ pour la sélection plate, avec $\approx\!48\%$ de jetons en moins ;
les régions sont connectées en interne à $95{,}5\%$.)
\end{abstract}

\section{Introduction}

Un agent de code à base de LLM opérant sur un dépôt doit décider de façon répétée
\emph{quoi mettre dans la fenêtre de contexte}. Le motif dominant est la
récupération : encoder des fragments de code, les classer par rapport à la tâche, et
concaténer le top-$k$ jusqu'à épuisement d'un budget de
jetons~\cite{lewis2020rag}. Cela traite le contexte comme une liste classée
unidimensionnelle. Deux modes de défaillance s'ensuivent. D'abord, la
\emph{dilution} : à mesure que les tâches s'allongent, la fenêtre accumule du code
globalement saillant mais sans rapport avec la tâche, et les modèles prêtent mal
attention au milieu des longs contextes~\cite{liu2023lost}. Ensuite,
l'\emph{incohérence} : les $k$ fragments les mieux classés ne forment pas
nécessairement une unité connexe de raisonnement --- une fonction peut être paginée
sans ses appelants, ses sites d'erreur ou les données dont elle dépend.

Les systèmes d'exploitation ont affronté le problème analogue il y a des décennies
et ont répondu par le \emph{working set} : l'ensemble des pages récemment
référencées qui doivent rester résidentes~\cite{denning1968working}. CCOS adopte
l'analogie pour le contexte LLM~\cite{packer2023memgpt} : le code source est analysé
en un graphe causal, les nœuds sont scorés, et une « fenêtre de contexte » bornée est
paginée comme la RAM~$\leftrightarrow$~VRAM. Contrairement à une pile de
récupération, chaque transition est enregistrée dans un journal d'événements chaîné
par hachage, si bien que la mémoire n'est pas une boîte noire mais un artefact
\emph{auditable et rejouable}. Nos contributions sont :

\begin{enumerate}
  \item Un \textbf{modèle formel et déterministe de régions de contexte}
  (\S\ref{sec:regions}, \S\ref{sec:determinism}) : la distance causale comme plus
  court chemin pondéré, l'appartenance à une région comme composante connexe du
  graphe des liens causaux inter-fichiers, et un théorème de déterminisme --- les
  régions et la fenêtre paginée sont une fonction pure du graphe, donc une session
  se reconstruit bit pour bit depuis son journal d'événements chaîné par hachage
  (rejeu à l'épreuve de la falsification).
  \item Des \textbf{sessions d'agent event-sourced avec débogage par voyage dans le
  temps} (\S\ref{sec:timetravel}) : chaque opération cognitive (ingestion, signal
  d'échec, rappel) est journalisée ; l'état exact du contexte à toute étape est
  reconstructible ; et un rappel peut être \emph{rejoué sous d'autres paramètres}
  pour demander si l'agent aurait mieux décidé --- la capacité qui manque
  structurellement à une pile de récupération probabiliste. À notre connaissance,
  c'est le premier traitement de l'\emph{assemblage du contexte lui-même} comme un
  sous-système rejouable et déboguable a posteriori --- un \emph{enregistreur de vol
  de l'attention} --- doté d'un \emph{point d'arrêt d'éviction} qui nomme l'étape et
  l'opération exactes où la vraie cause a été expulsée de la fenêtre budgétée.
  \item Un \textbf{harnais de validation honnête et un résultat négatif}
  (\S\ref{sec:protocol}) : sur $70$ commits réels de correction de bogues, la
  sélection causale \emph{ne bat pas} un récupérateur lexical TF-IDF pour placer les
  fichiers d'un correctif dans la fenêtre (elle fait jeu égal, et perd à budget
  serré), et un pivot par trace de crash est battu par un
  RAG-sur-le-message-d'erreur. Nous le rapportons clairement --- cela déplace la
  valeur de CCOS de la \emph{récupération} vers l'\emph{auditabilité}.
  \item Une \textbf{mesure de localité sans LLM} (\S\ref{sec:eval}) avec des chiffres
  réels et reproductibles, et un protocole falsifiable de condition
  \emph{suffisante} (Phase~4) que nous spécifions mais laissons ouvert.
\end{enumerate}

\section{Travaux connexes}\label{sec:related}

\paragraph{Génération augmentée par récupération.} Le RAG augmente un modèle
paramétrique d'un magasin non paramétrique et récupère des passages par
requête~\cite{lewis2020rag}. Self-RAG ajoute des jetons de réflexion qui contrôlent
la récupération et critiquent les générations~\cite{asai2023selfrag}. Ces approches
opèrent sur des fragments \emph{indépendants} ; la cohérence entre les éléments
récupérés n'est pas modélisée.

\paragraph{Récupération sensible aux graphes et à la structure.} GraphRAG construit
un graphe de connaissances d'entités et répond aux requêtes globales en résumant la
structure communautaire~\cite{edge2024graphrag}. Pour le code, les graphes de
propriétés unifient AST, flot de contrôle et flot de données~\cite{yamaguchi2014cpg}.
Les régions de CCOS s'inscrivent dans cette lignée mais visent la \emph{pagination} :
quelle sous-structure connexe rendre résidente pour une tâche, sous un budget de
jetons, avec une éviction déterministe.

\paragraph{Mémoire d'agent.} MemGPT présente le LLM comme un OS qui pagine entre une
fenêtre en contexte et un stockage externe~\cite{packer2023memgpt} ; les Generative
Agents récupèrent des souvenirs scorés par récence, importance et
pertinence~\cite{park2023generative} --- les mêmes facteurs que CCOS agrège en une
température de région. LangGraph~\cite{langgraph} structure les agents en graphes
d'étapes à état ; il orchestre le flot de \emph{contrôle}, là où CCOS structure la
\emph{mémoire}. Les deux sont complémentaires.

\paragraph{Gestion de la fenêtre de contexte.} L'éviction du cache KV maintient
résidents les jetons « gros frappeurs » ou puits~\cite{liu2023lost}, et
PagedAttention applique la pagination OS au cache KV~\cite{kwon2023paged}. Ces
techniques agissent au niveau du jeton à l'intérieur d'une passe avant ; CCOS agit au
niveau \emph{sémantique} à travers une session d'agent.

\section{Le substrat de contexte causal}\label{sec:substrate}

CCOS analyse le code source en un \emph{graphe de mémoire causale} dirigé $G=(V,E)$.
Un nœud $v\in V$ est un fichier, un module, un symbole ou une dépendance externe,
avec des champs scalaires $\mathrm{imp}(v)$ (importance de base),
$\mathrm{fail}(v)\in[0,1]$ (pertinence d'échec) et $\mathrm{rec}(v)\in[0,1]$
(récence). Une arête $e=(u\!\to\!w)\in E$ porte un poids $w(e)\in(0,1]$ et un type
(contenance, dépendance, référence, causation). Le noyau attribue à chaque nœud un
score causal
\begin{equation}
\mathrm{score}(v) = \mathrm{clamp}\big(0.15\,\mathrm{imp}(v) + 0.50\,\mathrm{fail}(v)
  + 0.30\,\mathrm{rec}(v) + 0.05\ln(1{+}\mathrm{acc}(v)),\,0,1\big),
\label{eq:score}
\end{equation}
où $\mathrm{acc}(v)$ est le compteur d'accès. Les fautes se propagent le long des
arêtes, $\mathrm{fail}$ décroît avec une horloge logique, et chaque transition d'état
est ajoutée à un journal d'événements chaîné par hachage permettant un rejeu
déterministe~\cite{haber1991timestamp}. Les identifiants de nœuds sont à espace de
noms (\code{file:p}, \code{mod:p:n}, \code{use:p:path}, \code{sym:p:n},
\code{dep:root}) ; le fichier propriétaire d'un nœud est récupérable depuis son
identifiant. CCOS v0.2 pagine les nœuds de plus haut $\mathrm{score}$ : une politique
plate et 1-D. Nous rendons désormais la sélection spatiale.

\section{Régions de contexte}\label{sec:regions}

\subsection{Distance causale}

\begin{definition}[Distance causale]
Soit $\hat G$ le multigraphe non orienté sur $V$ induit par $E$, et attribuons à
chaque arête le coût $c(e) = -\ln w(e) \ge 0$, de sorte qu'un lien causal plus fort
soit un pas plus court. La \emph{distance causale} $\dcausal(u,v)$ est le coût total
minimum sur tous les chemins $u$--$v$ dans $\hat G$, et $+\infty$ s'il n'en existe
aucun. La distance en sauts non pondérée $\mathrm{hops}(u,v)$ est définie de façon
analogue avec des coûts unitaires.
\end{definition}

$c(e)\ge 0$ puisque $w(e)\le 1$, donc $\dcausal$ est une véritable métrique de plus
court chemin sur chaque composante connexe (non négative, symétrique, inégalité
triangulaire). Le \emph{voisinage causal à $k$ sauts} d'une cible $t$ est
\begin{equation}
\Nk(t) = \{\, v \in V : \mathrm{hops}(t,v) \le k \,\}.
\end{equation}
$\Nk(t)$ est la vérité-terrain « ce qui est causalement pertinent pour une tâche en
$t$ » utilisée dans l'évaluation (\S\ref{sec:eval}).

\subsection{Appartenance à une région}

Les concentrateurs de dépendances externes (p. ex.\ \code{dep:std}) relient presque
tout et ne doivent pas effondrer le graphe en une seule composante. Nous les
séparons donc.

\begin{definition}[Lien causal inter-fichiers]
Pour un nœud non externe $v$, soit $\phi(v)$ son fichier propriétaire. Deux fichiers
$f,g$ sont \emph{directement liés}, noté $f \approx g$, si et seulement s'il existe
une arête $(u\!\to\!w)\in E$ avec $\phi(u)=f$, $\phi(w)=g$, $f\neq g$, et ni $u$ ni
$w$ n'est un nœud de dépendance externe.
\end{definition}

\begin{definition}[Région]\label{def:region}
Soit $\approx^{*}$ la clôture réflexive-transitive de $\approx$ sur l'ensemble des
clés de fichiers. Une \emph{région} est l'ensemble de tous les nœuds dont les
fichiers appartiennent à une même classe d'équivalence de $\approx^{*}$ ; tous les
nœuds de dépendance externe forment une région supplémentaire. Deux nœuds
appartiennent à la même région si et seulement si leurs fichiers sont dans la même
composante connexe du graphe des liens causaux inter-fichiers.
\end{definition}

\begin{proposition}[Les régions partitionnent les nœuds]
Les régions de la Définition~\ref{def:region} forment une partition de $V$ : chaque
nœud appartient à exactement une région.
\end{proposition}
\begin{proof}
$\approx^{*}$ est une relation d'équivalence sur les clés de fichiers, donc ses
classes partitionnent les clés ; le seau externe est une classe distincte par
construction. Associer chaque nœud non externe à la classe de son fichier et chaque
nœud externe à la classe externe est une fonction totale vers des blocs disjoints,
donc une partition. \qed
\end{proof}

Par défaut (uniquement des arêtes de contenance et d'importation), chaque fichier
source est sa propre région. Une véritable dépendance inter-fichiers ou un échec
propagé fusionne des fichiers en une seule région multi-fichiers --- la « zone de
connaissance » qu'un agent devrait réveiller ensemble.

\subsection{Scalaires de région}

Pour une région $R$ d'ensemble de membres $M$ :
\begin{align}
\mathrm{heat}(v)        &= 0.5\,\mathrm{score}(v) + 0.3\,\mathrm{fail}(v) + 0.2\,\mathrm{rec}(v), \\
\mathrm{temp}(R)        &= \mathrm{clamp}\!\Big(\tfrac{1}{|M|}\textstyle\sum_{v\in M}\mathrm{heat}(v),\,0,1\Big), \label{eq:temp}\\
\mathrm{dens}(R)        &= \frac{|\{\,e\in E : \text{les deux extrémités} \in M\,\}|}{|M|}. \label{eq:dens}
\end{align}
La température mesure à quel point une région est « éveillée » ; la densité est sa
cohésion causale interne (arêtes internes par membre). L'activation réchauffe une
région ($\mathrm{temp}\mathrel{+}=0.25$, plafonné) et enregistre un tic logique ; une
étape de refroidissement multiplie les températures par un facteur de décroissance et
évince les régions sous un plancher.

\subsection{Politique d'admission dynamique}\label{sec:policy}

Le seuil statique $0.6$ devient une fonction de la pression de jetons $u\in[0,1]$ (la
fraction utilisée du budget) et de la complexité de tâche $\kappa\in[0,1]$ :
\begin{align}
\theta(u,\kappa)       &= \mathrm{clamp}(0.6 + 0.3\,u - 0.2\,\kappa,\; 0.05,\; 0.95), \\
a(R)                   &= 0.55\,\mathrm{temp}(R) + 0.30\,\frac{\mathrm{dens}(R)}{1+\mathrm{dens}(R)} + 0.15\,\kappa, \\
\mathrm{admit}(R)      &\iff a(R) \ge \theta(u,\kappa).
\end{align}
Une région chaude et cohésive peut être admise même là où le $0.6$ statique la
rejetterait ; une fenêtre presque pleine élève $\theta$ pour que seules les régions
les plus chaudes entrent.

\section{Déterminisme et rejeu}\label{sec:determinism}

\begin{theorem}[Déterminisme régional]
La partition en régions, chaque scalaire de région des
équations~\eqref{eq:temp}--\eqref{eq:dens}, et l'historique d'activation sont des
fonctions pures, indépendantes de l'ordre, du graphe $G$ et de la suite d'événements
d'activation (logiquement horodatés). Par conséquent, étant donné le journal
d'événements d'une session, l'état du moteur se reconstruit à l'identique : si $G'$
est le graphe reconstruit depuis le journal et $L$ ses événements de région, alors
$\mathrm{replay\_from}(G',L)$ égale le moteur en direct qui a produit $L$.
\end{theorem}
\begin{proof}[Esquisse de preuve]
Le clustering énumère les nœuds et arêtes dans un ordre trié et calcule les
composantes connexes par un parcours en largeur trié, donc sa sortie est indépendante
de l'ordre d'itération du hachage. Les équations \eqref{eq:score}, \eqref{eq:temp} et
\eqref{eq:dens} sont une arithmétique déterministe sur les champs des nœuds et un
décompte d'arêtes qualifiantes. L'activation lit une horloge logique (un compteur),
jamais l'heure murale, et l'événement \code{RegionActivated} émis enregistre le tic
exact. La reconstruction re-clusterise $G'$ (état de base identique, puisque
$G'\!=\!G$ structurellement) et applique les activations enregistrées dans l'ordre du
journal ; chaque étape est la même fonction pure des mêmes entrées, donc l'état
résultant est identique. Le test d'intégration
\code{replay\_reconstructs\_identical\_engine} vérifie
$\mathrm{engine}=\mathrm{replay\_from}(G',L)$ par égalité structurelle. \qed
\end{proof}

Cela étend les garanties existantes de CCOS : le journal d'événements principal est
chaîné par hachage et à l'épreuve de la falsification, si bien que l'historique des
régions est auditable, et un test de 10\,000 cycles n'exhibe aucune dérive du nombre
de régions ni des températures.

\paragraph{Une frontière d'entrée auditable.} Le même substrat déterministe et
rejouable durcit le texte qu'un agent ingère. Les vecteurs d'injection par
caractères cachés --- surcharges bidirectionnelles (l'attaque \emph{Trojan
Source}, CVE-2021-42574), formatage de largeur nulle, et le bloc Unicode
\emph{Tags} servant à la contrebande d'ASCII invisible --- sont dé-obfusqués à
l'ingestion en littéraux explicites et visibles (\code{[U+202E RLO]}) plutôt que
silencieusement supprimés, et les \emph{findings} sont consignés dans le même
journal chaîné par hachage, de sorte qu'un rejeu reproduit l'état nettoyé. Un
classifieur linéaire en log-espace en aval --- la forme close du Naïve Bayes
multinomial, $\mathrm{logit}=b+W\!\cdot\!X$ sur un vecteur de caractéristiques
par \emph{hashing trick}, ses poids verrouillés dans un blob vérifié par somme de
contrôle --- ajoute un signal d'injection \emph{déterministe et décomposable
forensiquement} (un red-team en données réservées mesure $F_1=0{,}90$ ; précision
$0{,}87$, rappel $0{,}93$). Nous l'affirmons délibérément : c'est un \emph{signal,
pas une solution} --- un modèle sac-de-caractéristiques est contourné par
paraphrase, et aucune passe au niveau caractère ne traite l'injection sémantique ;
la vraie mitigation est la séparation des privilèges dans l'hôte. La contribution
n'est pas un détecteur inédit, mais que la dé-obfuscation et le score héritent de
la propriété distinctive de CCOS : toute décision de sécurité est reproductible et
auditable jusqu'à la caractéristique exacte qui l'a fait basculer.

\section{Sessions event-sourced et débogage par voyage dans le temps}\label{sec:timetravel}

Le déterminisme n'est pas seulement une propriété de correction ; il est le fondement
de la capacité qui, comme notre évaluation l'argumentera, distingue CCOS d'une pile
de récupération. Une \code{AgentSession} enregistre chaque opération cognitive qu'un
agent effectue sur sa mémoire --- \code{Ingest}, \code{SignalFailure}, \code{Recall}
--- comme une chronologie ordonnée. Parce que chaque opération est une fonction
déterministe de l'état antérieur, la mémoire après les $n$ premières opérations se
reconstruit exactement en rejouant celles qui mutent (\code{replay\_to(n)}) : on peut
\emph{rembobiner l'esprit de l'agent} à n'importe quelle étape.

L'opération qui compte pour le débogage est le contrefactuel.
\code{recall\_what\_if($n$, $q$, $b$)} rejoue la mémoire jusqu'à l'étape $n$ et
relance un rappel avec une requête $q$ ou un budget de jetons $b$ différents,
renvoyant la fenêtre que l'agent \emph{aurait} vue. Quand un agent émet un mauvais
correctif à l'étape 15, un opérateur peut rejouer son contexte exact à l'étape 14,
élargir le budget ou changer l'ancre, et observer si la décision s'améliore --- un
débogueur en boucle fermée et reproductible pour la mémoire de travail d'un agent.

\paragraph{Un point d'arrêt sur l'attention.} Le rejeu localise aussi la
\emph{dérive}. Le point d'observation \code{missing} parcourt la chronologie pour
trouver l'étape exacte où un nœud donné --- typiquement la vraie cause d'un échec ---
a été expulsé de la fenêtre budgétée par la pression concurrente, en nommant
l'opération déclenchante et l'écart de jetons ; une vue \code{energy} complémentaire
révèle la migration de chaleur causale à travers le graphe qu'un diff au niveau
fichier manque. C'est, de fait, un \emph{point d'arrêt sur l'attention d'un agent} :
l'instant précis où la bonne information a quitté la fenêtre. Une pile de récupération
opaque ne peut montrer que le contexte final, corrompu ; CCOS peut désigner l'étape
--- et l'opération --- où cela a déraillé.

Une pile RAG ou de framework mute son magasin de façon probabiliste au cours d'une
session et ne conserve aucun journal de transitions canonique et rejouable, de sorte
que la même question (\emph{pourquoi, et quand, la représentation s'est-elle
corrompue ?}) n'y a pas de réponse. Nous ne prétendons pas que cela améliore le
succès des tâches ; nous prétendons que c'est une propriété que les alternatives
n'ont pas, et que c'est le lieu honnête de la contribution de CCOS
(\S\ref{sec:protocol}).

\section{Implémentation}\label{sec:impl}

Le moteur fait $\approx\!600$ lignes de Rust léger en dépendances, en couche au-dessus
du graphe, sans modification du parseur, du garde, du constructeur incrémental ou du
cœur event-sourcing. La surface publique est : \code{ContextRegionEngine}
(\code{cluster\_nodes}, \code{initialize\_regions}, \code{activate\_region},
\code{tick\_cooldown}, \code{replay\_from}), les types de données \code{ContextRegion}
/ \code{ContextPoint} / \code{ContextWindow}, la \code{ContextPolicy}, et cinq
nouvelles variantes d'événement (\code{RegionCreated}, \code{RegionActivated},
\code{RegionMerged}, \code{RegionEvicted}, \code{ContextWindowGenerated}). Une CLI
\code{ccos regions} expose le clustering, l'activation et le rapport de localité.
L'ensemble de la crate se compile sans avertissement sous \code{clippy -D warnings}
et passe 364 tests.

\section{Évaluation : ce que nous pouvons mesurer aujourd'hui}\label{sec:eval}

Nous séparons les résultats \emph{mesurés} (cette section, entièrement reproductible
et sans LLM) des gains \emph{hypothétiques} au niveau de l'agent (\S\ref{sec:protocol}).

\paragraph{Dispositif.} Nous opérons sur la propre arborescence source de CCOS
($|V|=705$ nœuds, $|E|=822$ arêtes, donnant $26$ régions). Pour un nœud cible $t$,
nous prenons $\Nk(t)$ avec $k=2$ comme vérité-terrain et comparons deux stratégies de
sélection au budget de taille de la région : \textbf{plate} --- les $|R|$ meilleurs
nœuds par $\mathrm{score}$ (CCOS v0.2 / récupération classée classique) --- et
\textbf{région} --- les membres de la région de $t$. Nous rapportons la précision
causale $|S\cap\Nk|/|S|$, le rappel $|S\cap\Nk|/|\Nk|$, et les jetons que chaque
stratégie nécessite pour \emph{couvrir} $\Nk$. Tous les chiffres sont émis par
\code{scripts/region\_benchmark.sh} et sont déterministes.

\begin{table}[h]
\centering
\begin{tabular}{lcc}
\toprule
Stratégie & précision causale & rappel causal \\
\midrule
plate (v0.2 / classée) & 0.021 & 0.347 \\
\textbf{région (v0.3)} & \textbf{0.057} & \textbf{0.972} \\
\bottomrule
\end{tabular}
\caption{Localité sur \code{src/} ($k=2$, 12 cibles). La sélection par région couvre
$97\%$ du voisinage causal d'une tâche contre $35\%$ pour la sélection plate, avec
$\approx\!48\%$ de jetons en moins pour une couverture égale.}
\label{tab:locality}
\end{table}

\paragraph{Cohésion et coût.} Les régions atteignent une \textbf{densité causale
moyenne de $0{,}955$} (Éq.~\ref{eq:dens}, normalisée) --- elles sont presque
entièrement connectées en interne, c.-à-d.\ de véritables grappes causales plutôt que
des groupes de fichiers arbitraires. Construire la carte des régions pour $705$ nœuds
prend $\approx\!20$\,ms ($54{,}5$ constructions/s) ; une seule activation fait
$\approx\!20$\,ms.

\paragraph{Lecture honnête.} La précision absolue est faible ($0{,}06$) parce que le
parseur v0.2 n'émet que des arêtes de contenance et d'importation, donc $\Nk(t)$ est
minuscule (souvent $2$ à $3$ nœuds) tandis qu'une région couvre un fichier entier. Les
gains robustes et réels sont le \emph{rappel} ($0{,}97$ contre $0{,}35$) et
l'\emph{efficacité en jetons} ($\approx\!48\%$) : une région contient de façon fiable
le voisinage causal d'une tâche et paie moins de jetons pour le faire. Des arêtes
sémantiques plus riches (graphe d'appels / flot de données, item P1.3 de la feuille de
route) affineraient $\Nk$ et constituent le principal levier pour élever la précision ;
nous le signalons plutôt que de le cacher.

\section{Simulation d'hypothèse sous un oracle déclaré (sans LLM)}\label{sec:sim}

La thèse centrale --- \emph{la mémoire causale par régions aide-t-elle un agent sur de
longues tâches multi-fichiers ?} --- nécessite in fine des déroulements de LLM
(\S\ref{sec:protocol}). Mais sa \emph{condition nécessaire} est testable dès
maintenant, sans LLM : \textbf{un agent ne peut résoudre une tâche dont le contexte
causal requis est absent de sa fenêtre}. Nous mesurons, sous un oracle explicite, si
chaque stratégie de sélection \emph{place le contexte causal requis dans la fenêtre}.

\paragraph{Dispositif.} Nous générons des dépôts synthétiques modulaires : $S$
sous-systèmes indépendants, chacun un ensemble de fichiers reliés par une chaîne
causale inter-fichiers plus des symboles \emph{leurres} à haut score ; il n'y a pas
d'arêtes entre sous-systèmes, donc chaque sous-système est une région causale bornée.
La structure causale est \emph{découplée} de la structure lexicale --- un voisin de
chaîne ne partage aucun jeton d'identifiant avec la cible, seulement une arête ---
modélisant une dépendance causalement essentielle mais lexicalement dissemblable. Une
tâche de \emph{diamètre} $d$ requiert la fenêtre $\pm d$ le long de sa chaîne (jusqu'à
$2d{+}1$ fichiers) ; le budget contient $\approx$ un sous-système, bien moins que le
dépôt. L'\textbf{oracle} est $R(t)\subseteq S$.

Surtout, les pipelines de récupération localisent le code à partir d'une
\emph{requête} textuelle, tandis qu'une mémoire de type OS \emph{s'ancre} sur le
signal de l'espace de travail (le fichier actif, un test en échec). Nous modélisons les
deux : chaque tâche porte une requête (un sac de jetons) et une ancre (le nœud du
fichier actif), et nous exécutons deux scénarios --- \textbf{propre} (la requête pointe
vers la cible) et \textbf{bruité} (un leurre piège dans un sous-système sans rapport
surclasse lexicalement la cible). Six stratégies se partagent le budget :
\code{rag-dense} (top-$k$ lexical), \code{rag-hybrid} (lexical $+$ score causal),
\code{graphrag-1hop} (meilleur résultat $+$ un saut), \code{graphrag-bfs} (expansion
non bornée depuis le meilleur résultat), \code{ccos-from-query} (région CCOS du
meilleur résultat \emph{lexical} --- une ablation), et \code{ccos-region} (région CCOS
de l'\emph{ancre}). Amorcées et déterministes ; reproduire avec \code{ccos
experiment}.

\begin{table}[h]
\centering
\begin{tabular}{lcccccc}
\toprule
& \multicolumn{4}{c}{succès au diamètre $d$} & \multicolumn{2}{c}{global} \\
\cmidrule(lr){2-5}\cmidrule(lr){6-7}
Stratégie & $d{=}1$ & $d{=}2$ & $d{=}3$ & $d{=}4$ & succ. & couv. \\
\midrule
\multicolumn{7}{l}{\emph{Requête propre (pointe vers la cible)}}\\
\code{rag-dense}      & 0.00 & 0.00 & 0.00 & 0.00 & 0.00 & 0.19 \\
\code{rag-hybrid}     & 1.00 & 0.00 & 0.00 & 0.00 & 0.23 & 0.65 \\
\code{graphrag-1hop}  & 1.00 & 0.00 & 0.00 & 0.00 & 0.23 & 0.58 \\
\code{graphrag-bfs}   & 1.00 & 1.00 & 1.00 & 1.00 & 1.00 & 1.00 \\
\code{ccos-from-query}& 1.00 & 1.00 & 1.00 & 1.00 & 1.00 & 1.00 \\
\code{ccos-region}    & 1.00 & 1.00 & 1.00 & 1.00 & 1.00 & 1.00 \\
\midrule
\multicolumn{7}{l}{\emph{Requête bruitée (un leurre surclasse lexicalement la cible)}}\\
\code{rag-dense}      & 0.00 & 0.00 & 0.00 & 0.00 & 0.00 & 0.19 \\
\code{rag-hybrid}     & 1.00 & 0.00 & 0.00 & 0.00 & 0.23 & 0.65 \\
\code{graphrag-1hop}  & 0.00 & 0.00 & 0.00 & 0.00 & 0.00 & 0.19 \\
\code{graphrag-bfs}   & 0.00 & 0.00 & 0.00 & 0.00 & 0.00 & 0.00 \\
\code{ccos-from-query}& 0.00 & 0.00 & 0.00 & 0.00 & 0.00 & 0.00 \\
\textbf{\code{ccos-region}} & \textbf{1.00} & \textbf{1.00} & \textbf{1.00} & \textbf{1.00} & \textbf{1.00} & \textbf{1.00} \\
\bottomrule
\end{tabular}
\caption{Simulation d'hypothèse ($800$ tâches, graine $42$, budget $\approx$ un
sous-système). Succès $=$ l'ensemble causal requis $R(t)$ est dans la fenêtre. Sous
une requête propre, toute méthode sensible à la structure fait jeu égal ; sous une
requête trompeuse, seule la région ancrée sur l'espace de travail survit.}
\label{tab:sim}
\end{table}

\paragraph{Constats (et leurs limites honnêtes).}
\textbf{(1)} La récupération lexicale (\code{rag-dense}) \emph{échoue} sur les tâches
causales inter-fichiers ($0\%$ ; couverture $0{,}19$) : la similarité ne peut faire
émerger un contexte causalement essentiel mais lexicalement dissemblable --- la
prémisse de l'hypothèse. (Les victoires à $d{=}1$ de \code{rag-hybrid} viennent du
\emph{score causal} qu'il emprunte, non de la similarité lexicale, raison pour
laquelle elles persistent sous le bruit.)
\textbf{(2)} La valeur de la sélection sensible à la structure \emph{croît avec le
diamètre} : l'expansion à un saut ne résout que $d{=}1$, tandis que la pagination
causale complète (\code{graphrag-bfs}, \code{ccos-region}) résout tous les $d$ --- la
direction de H2.
\textbf{(3) Le jeu égal en cas propre.} Sous une requête propre, \code{graphrag-bfs},
\code{ccos-from-query} et \code{ccos-region} atteignent toutes $1{,}00$ : le levier est
la \emph{structure} causale, \textbf{pas CCOS en soi}.
\textbf{(4) La séparation en cas bruité.} Sous une requête \emph{trompeuse}, toute
méthode qui localise le code lexicalement s'effondre à $0\%$ --- y compris le robuste
\code{graphrag-bfs} \emph{et} l'ablation \code{ccos-from-query} (CCOS amorcé sur la
requête). Seule \code{ccos-region}, qui s'ancre sur le signal de l'espace de travail
plutôt que sur la requête, survit à $1{,}00$. L'ablation isole le différenciateur :
c'est la \emph{source de l'ancre} (un signal structurel d'espace de travail contre une
requête lexicale), non la machinerie des régions. C'est le régime réaliste --- une
description de tâche nomme rarement la cause distante exacte, et une mémoire de type OS
qui suit le working set actif est robuste là où la récupération au moment de la requête
est trompée.
\textbf{(5) Hypothèses, déclarées.} Cela crédite CCOS d'une ancre fiable (le fichier
actif / le test en échec), qu'une mémoire au niveau OS possède mais qu'un pur pipeline
de récupération n'a pas ; et cela suppose des dépôts \emph{modulaires} aux régions
séparables (densité $0{,}955$ sur du vrai code, \S\ref{sec:eval}) --- une région géante
monolithique dépassant le budget effondre l'avantage (observé, rapporté).
\textbf{(6)} Tout au long, c'est une \emph{simulation sous un oracle déclaré} : elle
teste la condition \emph{nécessaire} (récupération), non la condition
\emph{suffisante} (génération) de la \S\ref{sec:protocol}.

\section{Évaluation LLM réelle : une première mesure embarquée}\label{sec:protocol}

La simulation (\S\ref{sec:sim}) a établi la condition nécessaire. La condition
suffisante --- qu'un agent LLM \emph{résolve} ensuite plus de tâches avec une mémoire
par régions qu'avec une récupération par fragments --- requiert des déroulements de
modèle. Nous \textbf{implémentons le harnais} (\code{ccos eval}) et rapportons un
\textbf{premier essai LLM réel} contre un modèle local.

\paragraph{Harnais implémenté.} Chaque tâche est un minuscule projet multi-fichiers
encodant une \emph{chaîne causale arithmétique} : une constante de base dans un fichier
est transformée à travers une chaîne de fonctions d'une ligne réparties dans des
fichiers distincts, et la question « quel entier la dernière fonction renvoie-t-elle ? »
n'est répondable \emph{qu'en} lisant toute la chaîne (la cause distante comprise). La
notation est donc une correspondance exacte sur un entier --- objective et
automatisable, sans exécution de code. Les six mêmes stratégies de la \S\ref{sec:sim}
assemblent la fenêtre de contexte sous un budget de jetons (avec la division de requête
propre/bruitée), la fenêtre est envoyée à un modèle (tout point d'accès compatible
OpenAI, Anthropic-Messages ou Ollama), et nous enregistrons trois métriques par
diamètre : le \textbf{succès de tâche} (entier correct), la \textbf{couverture oracle}
(chaîne requise $\subseteq$ fenêtre --- indépendante du modèle), et
l'\textbf{hallucination de symbole} (la réponse cite une fonction absente du projet).

\paragraph{Dispositif.} Nous exécutons \code{ccos eval} en embarqué sur une NVIDIA
Jetson AGX Thor contre \code{qwen2.5:7b-instruct} servi par Ollama (température $0$),
avec $20$ tâches, graine $7$, un budget de $600$ jetons, des diamètres de $1$ à $4$,
dans les deux régimes propre et bruité. Le Tableau~\ref{tab:real} rapporte l'essai.

\begin{table}[h]
\centering
\begin{tabular}{lcccccr}
\toprule
& \multicolumn{4}{c}{succès au diamètre $d$} & & \\
\cmidrule(lr){2-5}
Stratégie & $d{=}1$ & $d{=}2$ & $d{=}3$ & $d{=}4$ & couv. & jet. \\
\midrule
\multicolumn{7}{l}{\emph{Requête propre (nomme la fonction cible)}}\\
\code{rag-dense}      & 0.12 & 0.00 & 0.00 & 0.00 & 0.00 & 519 \\
\code{rag-hybrid}     & 0.00 & 0.00 & 0.00 & 0.00 & 0.00 & 521 \\
\code{graphrag-1hop}  & 1.00 & 0.00 & 0.00 & 0.00 & 0.40 & 270 \\
\code{graphrag-bfs}   & 1.00 & 0.00 & 0.17 & 0.00 & 1.00 & 402 \\
\code{ccos-from-query}& 1.00 & 0.00 & 0.17 & 0.00 & 1.00 & 402 \\
\code{ccos-region}    & 1.00 & 0.00 & 0.17 & 0.00 & 1.00 & 402 \\
\midrule
\multicolumn{7}{l}{\emph{Requête bruitée (un leurre surclasse lexicalement la cible)}}\\
\code{rag-dense}      & 0.00 & 0.00 & 0.00 & 0.00 & 0.00 & 531 \\
\code{rag-hybrid}     & 0.00 & 0.00 & 0.00 & 0.00 & 0.00 & 521 \\
\code{graphrag-1hop}  & 0.00 & 0.00 & 0.00 & 0.00 & 0.00 & 165 \\
\code{graphrag-bfs}   & 0.00 & 0.00 & 0.00 & 0.00 & 0.00 & 165 \\
\code{ccos-from-query}& 0.00 & 0.00 & 0.00 & 0.00 & 0.00 & 165 \\
\textbf{\code{ccos-region}} & \textbf{1.00} & \textbf{0.00} & \textbf{0.17} & \textbf{0.00} & \textbf{1.00} & \textbf{402} \\
\bottomrule
\end{tabular}
\caption{Premier essai LLM réel : \code{qwen2.5:7b-instruct} via Ollama, en embarqué
(NVIDIA Jetson AGX Thor), $20$ tâches, graine $7$, budget $600$ jetons. « couv. » est
la couverture oracle indépendante du modèle ; « jet. » la moyenne des jetons d'entrée.
Le succès est une réponse entière en correspondance exacte.}
\label{tab:real}
\end{table}

\paragraph{Constats.}
\textbf{(1) La couverture se transfère de la simulation au texte réel.} La colonne de
couverture indépendante du modèle reproduit le Tableau~\ref{tab:sim} sur du contenu de
fichier réel : le RAG lexical couvre $0$, les méthodes sensibles à la structure couvrent
$1{,}00$ sous une requête propre, et sous le bruit \emph{seul} \code{ccos-region}
maintient $1{,}00$ --- la condition nécessaire, confirmée hors du simulateur.
\textbf{(2) Le succès est borné par le modèle, non par le contexte.} Là où la couverture
est $1{,}00$, le modèle de $7$B ne résout encore que $d{=}1$ ($1{,}00$) et une partie de
$d{=}3$ ($0{,}17$), manquant $d{=}2,4$ : avec toute la chaîne dans la fenêtre, l'échec
résiduel est le \emph{raisonnement arithmétique} --- un plafond de capacité du modèle,
non un échec de sélection. Surtout, aucune méthode ne réussit jamais \emph{sans}
couverture (succès $\subseteq$ couverture), ce qui est l'affirmation que le harnais est
bâti pour tester.
\textbf{(3) La séparation en cas bruité tient sur un vrai LLM.} Sous une requête
trompeuse, toute méthode pilotée par la requête s'effondre à $0\%$ de succès --- y
compris \code{graphrag-bfs} et l'ablation \code{ccos-from-query} --- tandis que
\code{ccos-region}, ancrée sur le signal de l'espace de travail, est la \emph{seule}
stratégie à succès non nul ($d{=}1$ à $1{,}00$). Les méthodes pilotées par la requête
paraissent même \emph{moins coûteuses} sous le bruit ($165$ contre $402$ jetons)
précisément parce qu'elles récupèrent avec assurance le mauvais fichier, plus petit ;
\code{ccos-region} dépense son budget sur la chaîne causalement correcte.
\textbf{(4) Limites honnêtes.} $20$ tâches ($\approx 5$ par diamètre) est un petit
échantillon, et un modèle de $7$B met au plancher la condition suffisante. Un modèle de
pointe (\code{deepseek-v4-pro}, en cours) ou un modèle local de $70$B devrait élever le
succès partout où la couverture est $1{,}00$, en \emph{affinant} --- non en changeant ---
la séparation, puisque la couverture fixe déjà le plafond.

\paragraph{Vers des benchmarks externes.} La suite à chaîne arithmétique isole la
question de la sélection ; les tests externes décisifs sont la résolution d'issues
SWE-bench~\cite{jimenez2023swebench} (un correctif passe les tests cachés) et une suite
contrôlée de \emph{bogues multi-fichiers} où le site de la faute, sa cause et son rayon
d'impact se trouvent dans des fichiers distincts. Les références partagent le LLM de
base et le budget de jetons : RAG classique~\cite{lewis2020rag} (top-$k$ cosinus sur
fragments), GraphRAG~\cite{edge2024graphrag} (résumés de communautés sur un graphe de
code), MemGPT~\cite{packer2023memgpt} (mémoire paginée de type OS), un agent
LangGraph~\cite{langgraph} avec un magasin vectoriel, et les régions CCOS. Métriques :
taux de résolution, jetons d'entrée jusqu'au succès, et un taux d'hallucination
d'ancrage de citation vérifié contre le graphe vérité-terrain. Nous avons commencé ceci
sur l'historique réel : \code{scripts/causal\_validation} extrait des commits de
correction, récupère l'arbre d'avant-correctif, injecte la faute, et score
$R_{\mathrm{cov}} = |F_{\mathrm{target}} \cap \mathrm{WorkingSet}_K| / |F_{\mathrm{target}}|$
--- la fraction des fichiers qu'un correctif a touchés que le working set borné
récupère. Sur l'historique de ce dépôt, le premier essai a exposé une limitation et sa
correction, toutes deux mesurées : avec une propagation d'échec uniquement en aval,
$R_{\mathrm{cov}}$ était plat à $0{,}33$ (seul le fichier de départ récupéré, car les
fichiers co-modifiés sont des importateurs \emph{en amont} atteints seulement par des
concentrateurs de dépendances). Résoudre les importations intra-crate en arêtes
fichier$\to$fichier et propager l'échec \emph{de façon bidirectionnelle} corrige cela.
Sur trois crates matures --- \code{fd}, \code{bat} et \code{hyperfine}, $70$ commits de
correction extraits --- l'effet est cohérent : à budget suffisant ($K{\ge}50$) les deux
changements élèvent $R_{\mathrm{cov}}$ à $0{,}85$--$1{,}0$ (depuis $0{,}50$--$0{,}84$ en
aval seul), tout en se \emph{diluant} à $0{,}19$--$0{,}28$ à un budget serré $K{=}20$.
\textbf{Mais face à la référence évidente, ce n'est pas une victoire.} En exécutant un
RAG lexical classique (cosinus TF-IDF sur le texte des fichiers, interrogé par le
fichier de la faute) au \emph{même budget de fichiers}, $R_{\mathrm{cov}}$ pour CCOS\,/\,RAG
est de $0{,}92/0{,}94$, $1{,}00/0{,}98$, $0{,}87/0{,}92$ à $K{=}50$ et fait jeu égal à
$K{=}100$, tandis qu'à $K{=}20$ le RAG est nettement devant ($0{,}20/0{,}56$ sur
\code{bat}, $0{,}20/0{,}73$ sur \code{hyperfine}). La sélection causale n'a donc
\emph{aucun avantage net de couverture} sur la similarité lexicale ici, et est pire à
budget serré : sur des bogues réels, les fichiers d'un correctif se ressemblent
lexicalement, si bien que le TF-IDF les récupère aussi --- la prémisse de la
\S\ref{sec:sim} d'une cause lexicalement dissemblable de son symptôme ne se reproduit pas
sur ces dépôts. La lecture honnête est que la condition \emph{nécessaire} tient pour la
grande majorité mais n'est \emph{pas} spécifique à CCOS ; un avantage réel devrait venir
des régimes que ce dispositif ne teste pas --- une requête dégradée/absente
(\S\ref{sec:sim}), un symptôme (test en échec) lexicalement loin de sa cause (nécessitant
le vrai amorçage piloté par le test, non l'heuristique du plus haut degré), ou la
condition \emph{suffisante} (Phase 4). Les espaces de travail multi-crates, désormais
liables, restent aussi à mesurer à grande échelle.

\paragraph{Condition suffisante (Phase 4) : la résolution fait jeu égal, l'efficacité
non.} Nous exécutons la moitié génération sur de vraies corrections mono-fichier de
\code{fd} : un agent reçoit un contexte à budget égal bâti de deux façons --- la région
causale de CCOS contre un RAG lexical top-$k$ --- on lui demande de réécrire le fichier
bogué, et il est noté selon que \code{cargo test} passe, avec une nouvelle tentative
\emph{compilateur dans la boucle} (à l'échec, une \emph{faute de page de contexte} analyse
l'erreur avec la couche de trace, injecte de la pression sur les fichiers fautifs, et
réamorce avec une fenêtre rafraîchie). Avec un faible modèle de $7$B et un seul essai, les
deux résolvent $0/15$ ; avec \code{qwen3-coder-30B} et trois tentatives, la boucle élève
les deux à $2/10$ ($20\%$) --- c'est le retour par faute de page, non le récupérateur, qui
débloque la résolution, et CCOS ne montre \textbf{aucun avantage de résolution} sur le RAG,
cohérent avec le résultat de récupération. L'\emph{efficacité} les sépare : à résolution
égale, la région auto-bornée de CCOS fait en moyenne $776$ jetons de contexte contre les
$5366$ qui remplissent le budget pour le récupérateur lexical --- une réduction de
$\mathbf{6{,}9\times}$. La même chose tient à l'échelle sans modèle : sur $51$ scénarios de
correction mono-fichier issus de \code{fd}, \code{bat} et \code{hyperfine}, CCOS assemble
$700$ à $1600$ jetons de contexte contre $\approx\!6000$ pour le RAG, une réduction de
$\mathbf{4{,}1}$ à $\mathbf{9{,}1\times}$. C'est le seul axe sur lequel CCOS domine : non
pas \emph{ce qu'}il récupère (une référence l'égale) mais \emph{combien peu} il lui faut,
parce que la région causale s'arrête au working set au lieu de remplir un top-$k$ jusqu'au
budget. Nous énonçons les réserves : l'échantillon de \emph{résolution} est minuscule
($n{=}10$, deux passes), et la référence remplit le budget par construction (un RAG à $k$
soigneusement réglé serait aussi plus clairsemé). L'affirmation défendable est plus étroite
mais réelle : CCOS \emph{s'auto-calibre} --- il se borne à la région causale sans $k$ ni
budget à régler --- donc il ne gaspille jamais la fenêtre, ce qui est précisément le but de
la pagination à la demande.

\paragraph{Hypothèses.}
\textbf{H1} (efficacité) : CCOS atteint un succès de tâche égal avec moins de jetons
d'entrée que RAG/GraphRAG, parce qu'une région couvre le voisinage causal avec un meilleur
rappel (Tableau~\ref{tab:locality}).
\textbf{H2} (succès à long horizon) : sur les tâches multi-fichiers, le taux de résolution
de CCOS dépasse le RAG par fragments d'une marge qui croît avec le diamètre causal de la
tâche.
\textbf{H3} (ancrage) : CCOS abaisse le taux d'hallucination de symbole, parce que la
fenêtre admise est un sous-graphe connexe de nœuds \emph{réels}.

\paragraph{Menaces à la validité.} La qualité d'une région est bornée par la qualité des
arêtes (\S\ref{sec:eval}) ; les résultats confondraient la \emph{politique} de région avec
le \emph{graphe} sous-jacent. Les ablations doivent donc fixer le graphe et ne varier que la
stratégie de sélection, et rapporter le diamètre causal par tâche pour que les gains soient
attribués à la cohérence, non au fait de récupérer plus de texte. L'ablation
\code{ccos-from-query} du Tableau~\ref{tab:real} est exactement ce contrôle --- même graphe
et même budget, seule l'ancre échangée --- et elle s'effondre avec les références lexicales
sous le bruit. Un résultat positif sur H1--H3 constituerait la contribution de recherche ;
un résultat nul validerait quand même l'infrastructure déterministe et auditable.

\section{Limites}

Nous sommes délibérément explicites sur ce qui n'est \emph{pas} montré.
\textbf{(1) Aucun avantage de couverture sur le RAG pour des bogues réels.} La métrique
phare sur données réelles ($R_{\mathrm{cov}}$, \S\ref{sec:protocol}) fait jeu égal avec un
simple récupérateur lexical TF-IDF à budget suffisant et lui est inférieure à budget serré ;
sur de vrais commits de correction, les fichiers qu'un correctif touche se ressemblent
lexicalement, donc la prémisse de la simulation d'une cause lexicalement dissemblable ne
tient pas. Les chiffres absolus forts sont une condition \emph{nécessaire} qu'une référence
satisfait aussi, non une victoire de CCOS.
(Une mesure complémentaire sur le code source du moteur lui-même isole une relation
différente, définie structurellement : un récupérateur lexical TF-IDF ne retrouve que
$49\%$ (rappel@$10$) des dépendances d'\emph{import} résolues par l'AST --- des paires
qui ne partagent souvent aucun vocabulaire, cosinus de dépendance $0{,}53$ contre $0{,}43$
aléatoire --- de sorte que l'angle mort du RAG est réel, mais il se situe hors de l'axe de
\emph{cohésion} d'un correctif que mesure la métrique phare, où les fichiers d'un correctif
partagent bien leur vocabulaire. La couche structurelle retrouve ces arêtes par construction ;
l'AST désormais par défaut est ce qui rend cette récupération complète.)
\textbf{(2) Nécessaire $\neq$ suffisant.} Nous ne mesurons jamais si une région \emph{aide
un agent à corriger} un bogue (Phase 4 : un correctif qui passe les tests cachés) --- seulement
si les fichiers pertinents sont récupérables.
\textbf{(3) Heuristique d'amorçage.} Le harnais injecte la faute dans le fichier modifié de
plus haut degré sortant, non dans le test réellement en échec, donc il n'exerce pas le régime
où un symptôme est lexicalement loin de sa cause --- le régime le plus favorable à la
sélection causale.
\textbf{(4) Échelle et statistiques.} Trois dépôts mono-crate, $n\!\approx\!20$ chacun ; les
différences sont rapportées avec leur écart-type mais non validées de façon croisée.
\textbf{(5) Moteur.} Le parseur utilise désormais \emph{par défaut} un vrai AST \texttt{syn}
--- mesuré $36{,}5\%$ plus précis que l'ancienne heuristique ligne sur le code source du moteur
lui-même (deux tiers contre la totalité du rappel d'imports), l'heuristique n'étant conservée
que comme repli pour les entrées non-Rust --- et l'ingestion a été durcie pour se reconstruire
bit à bit à l'identique entre processus (résolution d'imports triée ; centralité invariante à
l'ordre). Les régions restent à granularité fichier et ne fusionnent que sur des arêtes
inter-fichiers explicites ; les poids de score sont réglés à la main, bien que des raffinements
optionnels (désactivés par défaut) de centralité par vecteur propre et d'états de cycle de vie
des nœuds (\texttt{Stable}/\texttt{Working}/\texttt{Orphan}) soient désormais disponibles. Rien
de tout cela n'affecte les propriétés \emph{prouvées} (partition, déterminisme, rejeu) ni la
localité mesurée --- mais le dossier empirique d'un bénéfice pour l'agent en aval est, à ce
stade, \emph{non prouvé}.

\section{Conclusion}

Nous étions partis pour montrer qu'organiser la mémoire d'un agent par \emph{régions
causales} récupère le contexte à long horizon mieux que la récupération par fragments, avons
donné à la construction une définition précise et déterministe, et prouvé qu'elle se
reconstruit bit pour bit au rejeu. Nous avons ensuite testé l'affirmation de récupération
contre la référence évidente sur données réelles --- et elle n'a pas tenu : sur $70$ commits
réels de correction de bogues, un simple récupérateur lexical TF-IDF fait jeu égal avec la
sélection causale et la bat à budget serré, et un pivot par trace de crash perd contre un
RAG-sur-le-message-d'erreur, parce que sur du vrai code les fichiers d'un correctif et ses
messages d'erreur partagent leur vocabulaire. Nous rapportons le résultat négatif plutôt que
de l'enterrer. Ce qui survit n'est pas un meilleur récupérateur mais un type d'objet
différent : une mémoire de travail \emph{déterministe, rejouable, auditable} dans laquelle
l'état exact du contexte d'un agent peut être rembobiné et rejoué sous d'autres paramètres
--- un débogage par voyage dans le temps qu'une pile de récupération probabiliste ne peut
offrir. Que cette auditabilité, ou un avantage de contexte structuré au moment de la
\emph{génération} (Phase 4), produise un gain mesurable en aval reste ouvert. Nous publions
le moteur, le harnais de validation honnête et cet article afin que la distinction entre ce
qui est prouvé, ce qui est mesuré et ce qui est réfuté puisse être vérifiée plutôt
qu'affirmée.

\begin{thebibliography}{99}
\bibitem{lewis2020rag} P.~Lewis et al. Retrieval-Augmented Generation for
Knowledge-Intensive NLP Tasks. \emph{NeurIPS}, 2020. arXiv:2005.11401.
\bibitem{asai2023selfrag} A.~Asai et al. Self-RAG: Learning to Retrieve, Generate
and Critique through Self-Reflection. \emph{ICLR}, 2024. arXiv:2310.11511.
\bibitem{edge2024graphrag} D.~Edge et al. From Local to Global: A Graph RAG
Approach to Query-Focused Summarization. 2024. arXiv:2404.16130.
\bibitem{yamaguchi2014cpg} F.~Yamaguchi et al. Modeling and Discovering
Vulnerabilities with Code Property Graphs. \emph{IEEE S\&P}, 2014.
\bibitem{packer2023memgpt} C.~Packer et al. MemGPT: Towards LLMs as Operating
Systems. 2023. arXiv:2310.08560.
\bibitem{park2023generative} J.~S.~Park et al. Generative Agents: Interactive
Simulacra of Human Behavior. \emph{UIST}, 2023. arXiv:2304.03442.
\bibitem{langgraph} LangChain. LangGraph: Building Stateful, Multi-Actor
Applications with LLMs. Framework logiciel, 2023--2024.
\url{https://github.com/langchain-ai/langgraph}.
\bibitem{denning1968working} P.~J.~Denning. The Working Set Model for Program
Behavior. \emph{Communications of the ACM}, 11(5), 1968.
\bibitem{liu2023lost} N.~F.~Liu et al. Lost in the Middle: How Language Models Use
Long Contexts. \emph{TACL}, 2024. arXiv:2307.03172.
\bibitem{kwon2023paged} W.~Kwon et al. Efficient Memory Management for Large
Language Model Serving with PagedAttention. \emph{SOSP}, 2023. arXiv:2309.06180.
\bibitem{haber1991timestamp} S.~Haber and W.~S.~Stornetta. How to Time-Stamp a
Digital Document. \emph{Journal of Cryptology}, 3(2), 1991.
\bibitem{jimenez2023swebench} C.~E.~Jimenez et al. SWE-bench: Can Language Models
Resolve Real-World GitHub Issues? \emph{ICLR}, 2024. arXiv:2310.06770.
\end{thebibliography}

\end{document}
